\chapter{Dise\~no y Arquitectura}\label{cap_diseno}

\section{Arquitectura de alto nivel}

La arquitectura adoptada es cliente-servidor desacoplada, con un \textbf{frontend SPA} en React + TypeScript para presentaci\'on y experiencia de usuario, un \textbf{backend REST} en Spring Boot para l\'ogica de negocio y seguridad, y \textbf{persistencia relacional} con JPA/Hibernate y migraciones Flyway. Esta separaci\'on delimita responsabilidades y facilita evoluci\'on modular.

En backend se aplica dise\~no por capas
(controller, service, repository y entity),
que mejora mantenibilidad, testabilidad y capacidad de evoluci\'on.

\section{Vista de despliegue}

La vista de despliegue se organiza en bloques principales que cubren el \textbf{cliente web} (navegador del usuario final), el \textbf{frontend} como SPA est\'atica servida desde plataforma PaaS (Render), el \textbf{backend} como API REST desplegada en contenedor Docker en Render, los \textbf{datos} en PostgreSQL gestionada accesible solo desde backend, y los \textbf{servicios externos} de almacenamiento cloud, OneDrive y correo transaccional por API HTTP. Esta disposici\'on mantiene una frontera de seguridad clara entre capas.

La separaci\'on de componentes permite escalar y evolucionar cada capa
sin romper el resto del sistema,
manteniendo adem\'as una frontera de seguridad clara entre cliente,
l\'ogica de negocio y persistencia.

\begin{table}[!htbp]
\centering
\begin{tabular}{|P{2.9cm}|P{3.1cm}|P{5.6cm}|}
\hline
\textbf{Nodo} & \textbf{Tecnolog\'ia} & \textbf{Responsabilidad principal} \\
\hline
Cliente & Navegador moderno & Interacci\'on de usuario y consumo de API sobre HTTPS. \\
\hline
Frontend & React + Vite en Render & Renderizado de interfaz, estado de sesi\'on y experiencia por rol. \\
\hline
Backend & Spring Boot (Java 21) en Render & Exposici\'on de endpoints, reglas de negocio, seguridad y orquestaci\'on de persistencia. \\
\hline
Persistencia & Render PostgreSQL & Almacenamiento transaccional y trazabilidad de datos acad\'emicos. \\
\hline
Correo transaccional & Brevo API HTTP & Recuperaci\'on de contrase\~na y notificaciones sin depender de puertos SMTP salientes. \\
\hline
\end{tabular}
\caption{Vista de despliegue simplificada del sistema}
\label{tab:vista-despliegue}
\end{table}

\section{Dise\~no de dominio y datos}

El dominio se modela alrededor de la relaci\'on docente-estudiante-actividad, con \texttt{Usuario} como identidad base, \texttt{Profesor} y \texttt{Estudiante} como especializaciones operativas, \texttt{Curso}, \texttt{Grupo} y \texttt{Actividad} para el contexto acad\'emico, \texttt{Entregable} y \texttt{Entrega} para el ciclo de entrega, y \texttt{Feedback}, \texttt{Material} y entidades de soporte (tokens, notificaciones). Esta estructura facilita trazabilidad y coherencia de reglas de negocio.

La persistencia evoluciona de forma versionada mediante migraciones,
lo que permite incorporar nuevas capacidades sin perder consistencia hist\'orica.

\section{Dise\~no de API}

La API sigue orientaci\'on REST y separa contratos mediante DTOs. Se definen m\'odulos de endpoint coherentes por contexto de negocio, incluyendo \texttt{/api/auth} para login, refresh y recuperaci\'on de contrase\~na, \texttt{/api/actividades} para gesti\'on de actividades y consulta por usuario, \texttt{/api/entregables} para administraci\'on de entregables, y \texttt{/api/entregas} para subida, revisi\'on, evaluaci\'on y descarga de entregas. Esta segmentaci\'on facilita evoluci\'on de contratos por dominio.

El uso de DTOs evita exponer directamente entidades de persistencia,
reduce acoplamiento y simplifica evoluciones de contrato.

\section{Dise\~no de seguridad}

El dise\~no de seguridad incorpora mecanismos complementarios: autenticaci\'on basada en JWT (token de acceso y refresh token), filtros de validaci\'on en la cadena de seguridad, autorizaci\'on por rol y contexto acad\'emico, cifrado de contrase\~nas con BCrypt y limitaci\'on de tasa en operaciones sensibles de autenticaci\'on. Esta combinaci\'on reduce superficie de ataque sin penalizar la experiencia de uso.

Se configura CORS de forma expl\'icita y se diferencia entre rutas p\'ublicas y protegidas,
minimizando superficie de ataque sin comprometer usabilidad.

\section{Dise\~no del frontend}

El frontend se organiza en m\'odulos \texttt{pages}, \texttt{components}, \texttt{context} y \texttt{services} con objetivos de experiencia coherente para profesor y estudiante, protecci\'on de rutas seg\'un sesi\'on y rol, gesti\'on centralizada de autenticaci\'on y manejo uniforme de errores y estados de carga. Esta organizaci\'on evita duplicidades y mejora mantenibilidad.

La capa de red encapsula llamadas HTTP y gesti\'on de tokens,
incluyendo renovaci\'on autom\'atica de sesi\'on al expirar el token de acceso.

\FloatBarrier

Esta granularidad no es habitual en soluciones generalistas,
que suelen tratar todo como subida de fichero sin reglas contextuales.

\subsubsection*{Versionado y reenv\'io con trazabilidad}

Cada entrega queda registrada con versi\'on y fecha, manteniendo un
historial completo y una versi\'on activa por entregable.
Cuando el reenv\'io est\'a permitido, el sistema conserva la trazabilidad
sin sobrescribir el historial previo, facilitando auditor\'ia docente.

\subsubsection*{Operativa docente optimizada}

Se incorporan descargas masivas por entregable o actividad,
previsualizaci\'on de archivos en l\'inea y listado de contenido
de ZIP, reduciendo tareas manuales de comprobaci\'on y clasificaci\'on.

\subsubsection*{Robustez en el flujo de estudiante}

El flujo de entrega incorpora borradores autom\'aticos en cliente,
revalidaci\'on antes de enviar y mensajes de error contextualizados.
Esto minimiza p\'erdida de trabajo en conexiones inestables
y mejora la experiencia del alumnado.

\FloatBarrier
\section{Alternativas de arquitectura y stack}

\subsection{Arquitectura de aplicaci\'on}

Se valoraron dos enfoques principales: \textbf{A1 - SPA + API desacoplada} (opci\'on elegida) y \textbf{A2 - Aplicaci\'on monol\'itica server-side renderizada}. Esta comparativa permite sopesar flexibilidad de evoluci\'on frente a simplicidad inicial.

\textbf{A1} permite separar claramente la capa de experiencia de usuario
de la l\'ogica de negocio y persistencia, facilitando pruebas independientes,
despliegues diferenciados y evoluci\'on incremental.

\textbf{A2} reduce complejidad inicial de infraestructura,
pero limita flexibilidad de frontend y complica evoluci\'on de UX avanzada
en escenarios con roles y paneles diferenciados.

\subsection{Backend framework}

Se compararon principalmente \textbf{B1 - Spring Boot (Java)} (elegida), \textbf{B2 - Node.js (Express/Nest)} y \textbf{B3 - Django (Python)}. El an\'alisis consider\'o compatibilidad con seguridad, persistencia y tiempos de entrega.

La elecci\'on de \textbf{Spring Boot} se justifica por su ecosistema maduro para seguridad, JPA, validaci\'on y testing, por la integraci\'on natural con migraciones Flyway y por el buen soporte a arquitectura por capas y control transaccional. Estos factores reducen incertidumbre t\'ecnica en el alcance temporal disponible.

Node.js y Django son opciones v\'alidas,
pero para este proyecto requer\'ian m\'as decisiones de ensamblaje
en piezas de seguridad/persistencia para alcanzar un nivel equivalente
en el tiempo disponible.

\subsection{Frontend framework}

Alternativas consideradas: \textbf{F1 - React + TypeScript + Vite} (elegida), \textbf{F2 - Angular} y \textbf{F3 - Vue}. La selecci\'on se orient\'o a equilibrio entre velocidad de desarrollo y control de complejidad.

React con TypeScript y Vite ofrece buen equilibrio entre velocidad de desarrollo,
ecosistema y control de complejidad para una SPA de tama\~no medio.
Angular aporta una estructura m\'as opinionada,
pero con curva de entrada superior para el alcance temporal del TFG.
Vue mantiene una DX excelente, aunque en este caso la experiencia previa del equipo
y la alineaci\'on con librer\'ias de testing/CI favorecieron React.

\FloatBarrier
\section{Alternativas de persistencia y almacenamiento}

\subsection{Base de datos relacional}

Opciones comparadas: \textbf{D1 - PostgreSQL gestionado} (elegida en producci\'on), \textbf{D2 - MySQL gestionado} y \textbf{D3 - H2 embebida} (solo desarrollo/pruebas locales). La selecci\'on equilibra estabilidad productiva con agilidad local.

PostgreSQL se mantuvo como opci\'on principal por estabilidad, compatibilidad
con el stack y buen comportamiento en entornos PaaS.
H2 se reserva para contextos de desarrollo y pruebas, no para producci\'on.

Dentro de las alternativas de base de datos desplegadas,
inicialmente se decidi\'o utilizar \textbf{Neon} como proveedor PostgreSQL
por sus mejores caracter\'isticas percibidas para el arranque del proyecto:
aprovisionamiento r\'apido, enfoque serverless y facilidad de conexi\'on
desde un backend desplegado en PaaS.
No obstante, durante la validaci\'on operativa se detectaron limitaciones
en su plan gratuito que condicionaban la continuidad del despliegue acad\'emico.
Por ello se migr\'o a \textbf{Render PostgreSQL},
manteniendo PostgreSQL como decisi\'on t\'ecnica de base y cambiando solo
el proveedor gestionado.
La migraci\'on se realiz\'o sin cambios de c\'odigo de dominio,
apoy\'andose en variables de entorno y configuraci\'on portable del datasource,
lo que confirma bajo acoplamiento a proveedor concreto.

\subsection{Estrategia de esquema}

Se compararon \textbf{E1 - Flyway como autoridad de esquema} (elegida) y \textbf{E2 - Hibernate \texttt{ddl-auto=update} en producci\'on}. La decisi\'on prioriza trazabilidad y control expl\'icito del cambio.

Se adopta \textbf{E1} por trazabilidad y control expl\'icito de cambios,
reduciendo riesgo de divergencias de esquema entre entornos.

\subsection{Almacenamiento de archivos}

Alternativas evaluadas: \textbf{S1 - Almacenamiento local en servidor}, \textbf{S2 - Cloudinary} (elegida para producci\'on persistente) y \textbf{S3 - OneDrive/Microsoft Graph} (integraci\'on opcional por contexto acad\'emico). La opci\'on final combina persistencia en producci\'on con flexibilidad institucional.

La opci\'on local es simple y adecuada en desarrollo,
pero en plataformas con sistema de archivos ef\'imero (como despliegues free)
puede provocar p\'erdida de datos tras reinicio.
Por ello se adopta \textbf{Cloudinary} en producci\'on.

La integraci\'on \textbf{OneDrive} aporta valor en casos docentes
que requieren sincronizaci\'on con ecosistemas Microsoft,
pero se mantiene como capacidad configurable y no obligatoria del n\'ucleo.

\FloatBarrier
\section{Alternativas de seguridad y sesi\'on}

\subsection{Modelo de autenticaci\'on}

Se compararon dos enfoques principales: \textbf{A1 - JWT (access + refresh)} (elegido) y \textbf{A2 - Sesi\'on server-side tradicional}. La elecci\'on favorece coherencia con arquitectura SPA y escalado stateless.

En una SPA desacoplada, JWT con refresh ofrece una experiencia m\'as natural
para cliente API-first y facilita escalado stateless del backend.
La sesi\'on server-side reduce cierta complejidad de token,
pero exige estrategia adicional de estado distribuido en despliegues horizontales.

\subsection{Extensiones de seguridad}

Se han considerado funcionalidades avanzadas (inicio de sesi\'on federado con OAuth2/OIDC y 2FA)
como l\'inea de evoluci\'on.
La decisi\'on de no imponerlas en el MVP responde al criterio
de maximizar cobertura funcional core sin comprometer plazos.

\FloatBarrier
\section{Alternativas investigadas y descartadas}

\subsection{SharePoint para OneDrive corporativo}

Se valor\'o el uso de SharePoint como intermediario para acceder a cuentas
corporativas de OneDrive. La opci\'on facilitaba el encaje institucional,
pero implicaba perder capacidades clave que s\'i ofrece la integraci\'on directa,
como la selecci\'on de carpeta de destino o la gesti\'on flexible de rutas
durante la subida de entregas. Dado que estas funciones son parte del valor
diferencial de la plataforma, se descart\'o esta v\'ia.

\subsection{Registro de usuarios con OAuth}

Se estudi\'o habilitar registro de usuarios mediante OAuth/OIDC
para que el alumnado pudiera crear cuentas de forma aut\'onoma.
Finalmente se prioriz\'o el flujo habitual en entornos acad\'emicos,
donde el administrador crea usuarios y los asigna a cursos y grupos.
Esto evita altas desalineadas con la matr\'icula real y simplifica el control
de permisos y pertenencia a grupos.

\FloatBarrier
\section{Alternativas de despliegue y operaci\'on}

\subsection{Estrategia de despliegue}

Opciones principales: \textbf{P1 - PaaS con automatizaci\'on GitHub Actions + Render} (elegida), \textbf{P2 - VPS autogestionado con pipeline manual o Jenkins} y \textbf{P3 - Orquestaci\'on completa en Kubernetes}. La selecci\'on busca minimizar carga operativa y mantener viabilidad acad\'emica.

\textbf{P1} reduce carga operativa,
permite despliegue r\'apido y mantiene un coste ajustado para contexto acad\'emico.
La opci\'on elegida se concreta en Render como PaaS gestionado para el servicio
backend, el frontend est\'atico y PostgreSQL gestionado, manteniendo la
configuraci\'on por variables de entorno.
\textbf{P2} y \textbf{P3} ofrecen m\'as control,
pero su complejidad no aporta retorno proporcional dentro del alcance del TFG.

Dentro de esta estrategia, el correo transaccional se resuelve con Brevo
mediante API HTTP. Inicialmente se emple\'o SendGrid como proveedor
de correo, pero tras la finalizaci\'on de su plan gratuito se migr\'o
a Brevo. Se evalu\'o tambi\'en Resend, pero su plan gratuito requiere
un dominio verificado para enviar a destinatarios arbitrarios.
Esta alternativa se prefiere frente a SMTP en producci\'on
porque los PaaS pueden bloquear o restringir puertos propios del protocolo
de correo, mientras que las peticiones HTTPS salientes son compatibles con
el modelo operativo del despliegue.

\subsection{CI/CD y calidad}

Se compararon \textbf{C1 - CI/CD integrada en GitHub Actions} (elegida) y \textbf{C2 - CI externa (GitLab CI/Jenkins) con mayor administraci\'on}. La decisi\'on prioriza trazabilidad y menor sobrecarga operativa.

La alternativa elegida integra lint, tests, build, an\'alisis de seguridad,
mutaci\'on y publicaci\'on de artefactos en un flujo unificado y reproducible.

\FloatBarrier
\section{Comparativa multicriterio}

La Tabla~\ref{tab:comparativa-global} sintetiza la comparaci\'on global
entre enfoques de soluci\'on.

\begin{table}[!htbp]
\centering
\small
\begin{tabular}{|P{3.9cm}|c|c|c|}
\hline
\textbf{Criterio} & \textbf{Soluci\'on elegida} & \textbf{Alternativa intermedia} & \textbf{Alternativa avanzada} \\
\hline
Tiempo de entrega MVP & 5 & 3 & 2 \\
\hline
Mantenibilidad & 4 & 3 & 4 \\
\hline
Coste operativo & 5 & 3 & 2 \\
\hline
Escalabilidad potencial & 4 & 3 & 5 \\
\hline
Complejidad de operaci\'on & 4 & 3 & 1 \\
\hline
Adecuaci\'on al TFG & 5 & 3 & 2 \\
\hline
\end{tabular}
\caption{Comparativa global (escala 1-5, mayor es mejor)}
\label{tab:comparativa-global}
\end{table}

Interpretaci\'on:
la alternativa seleccionada no maximiza todos los ejes te\'oricos,
pero s\'i ofrece el mejor equilibrio entre valor entregado,
riesgo t\'ecnico y viabilidad temporal del proyecto.

\FloatBarrier
\section{Decisiones adoptadas y justificaci\'on final}

Las decisiones finales que mejor explican el resultado del proyecto son la arquitectura SPA + API desacoplada para separar responsabilidades, un backend Spring Boot con persistencia relacional y migraciones Flyway, un frontend React + TypeScript para productividad con control tipado, almacenamiento cloud persistente en producci\'on y local en desarrollo, y despliegue en PaaS con CI/CD automatizada y verificaciones de calidad. Esta combinaci\'on asegura equilibrio entre viabilidad temporal y robustez t\'ecnica.

Esta combinaci\'on ha permitido llegar a un sistema funcional,
con cobertura de pruebas alta y capacidad de evoluci\'on,
sin sobredimensionar infraestructura ni desviarse del objetivo acad\'emico.

\FloatBarrier
\section{Conclusiones del cap\'itulo}

La comparaci\'on confirma que la soluci\'on implementada no responde a una selecci\'on
arbitraria de herramientas,
sino a una evaluaci\'on razonada frente a alternativas reales.

El enfoque adoptado prioriza equilibrio entre robustez, coste y mantenibilidad,
lo que lo hace especialmente adecuado para un TFG de perfil profesional
con restricciones de tiempo y recursos.
	
	 


























